\documentclass[11pt, a4paper]{article}

% bfw-style.tex — gemeinsame Style-Definitionen für alle Handouts/Aufgabenhefte
% Voraussetzung: Kompilation mit xelatex (wegen fontspec)

\usepackage[ngerman]{babel}
\usepackage{geometry}
\geometry{a4paper, top=2.2cm, bottom=2.2cm, left=2.3cm, right=2.3cm}

\usepackage{fontspec}
% Fallback-Kette: Source Sans Pro -> Calibri -> Segoe UI -> Latin Modern Sans
% (Latin Modern ist immer mit MiKTeX vorhanden)
\IfFontExistsTF{Source Sans Pro}{%
  \setmainfont{Source Sans Pro}[Ligatures=TeX]%
}{\IfFontExistsTF{Calibri}{%
  \setmainfont{Calibri}[Ligatures=TeX]%
}{\IfFontExistsTF{Segoe UI}{%
  \setmainfont{Segoe UI}[Ligatures=TeX]%
}{%
  \setmainfont{Latin Modern Sans}[Ligatures=TeX]%
}}}
\IfFontExistsTF{Source Code Pro}{%
  \setmonofont{Source Code Pro}[Scale=0.92]%
}{\IfFontExistsTF{Consolas}{%
  \setmonofont{Consolas}[Scale=0.92]%
}{%
  \setmonofont{Latin Modern Mono}[Scale=0.92]%
}}

\usepackage{xcolor}
% Farbpalette: hell, freundlich, modern
\definecolor{bfwPrimary}{HTML}{0EA5A4}    % Petrol/Türkis - Primär
\definecolor{bfwPrimaryDark}{HTML}{0F766E} % dunkler Petrol für Headings
\definecolor{bfwAccent}{HTML}{F59E0B}     % Amber - Akzent
\definecolor{bfwSuccess}{HTML}{10B981}    % Grün - Erfolg / Lösung
\definecolor{bfwWarn}{HTML}{F97316}       % Orange - Achtung
\definecolor{bfwInk}{HTML}{0F172A}        % fast Schwarz
\definecolor{bfwInkSoft}{HTML}{475569}    % gedämpftes Grau
\definecolor{bfwBgSoft}{HTML}{F8FAFC}     % off-white
\definecolor{bfwBgTeal}{HTML}{ECFEFF}     % helles Türkis
\definecolor{bfwBgAmber}{HTML}{FEF3C7}    % helles Gelb
\definecolor{bfwBgRose}{HTML}{FFE4E6}     % helles Rosé für Eselsbrücken

\usepackage[colorlinks=true, linkcolor=bfwPrimaryDark, urlcolor=bfwPrimaryDark]{hyperref}
\usepackage{enumitem}
\setlist[itemize]{leftmargin=*, itemsep=2pt, topsep=2pt}
\setlist[enumerate]{leftmargin=*, itemsep=2pt, topsep=2pt}

\usepackage[most]{tcolorbox}
\usepackage{booktabs}
\usepackage{tikz}
\usetikzlibrary{decorations.pathmorphing, positioning, shapes.geometric, arrows.meta}

% Merkkästchen — wichtige Definition / Kernpunkt
\newtcolorbox{merksatz}[1][]{%
  enhanced, breakable,
  colback=bfwBgTeal,
  colframe=bfwPrimary,
  boxrule=0.6pt,
  arc=4pt,
  left=10pt, right=10pt, top=8pt, bottom=8pt,
  fonttitle=\bfseries\color{bfwPrimaryDark},
  title={\faLightbulb\ Merksatz},
  #1
}

% Eselsbrücke — Mnemoniks
\newtcolorbox{eselsbruecke}[1][]{%
  enhanced, breakable,
  colback=bfwBgRose,
  colframe=bfwAccent,
  boxrule=0.6pt,
  arc=4pt,
  left=10pt, right=10pt, top=8pt, bottom=8pt,
  fonttitle=\bfseries\color{bfwWarn},
  title={Eselsbrücke},
  #1
}

% Beispiel-Box
\newtcolorbox{beispiel}[1][]{%
  enhanced, breakable,
  colback=bfwBgSoft,
  colframe=bfwInkSoft,
  boxrule=0.4pt,
  arc=3pt,
  left=10pt, right=10pt, top=8pt, bottom=8pt,
  fonttitle=\bfseries\color{bfwInk},
  title={Beispiel},
  #1
}

% Lösungs-Box (für Aufgabenhefte)
\newtcolorbox{loesung}[1][]{%
  enhanced, breakable,
  colback=bfwBgAmber!40,
  colframe=bfwSuccess,
  boxrule=0.6pt,
  arc=3pt,
  left=10pt, right=10pt, top=8pt, bottom=8pt,
  fonttitle=\bfseries\color{bfwSuccess!70!black},
  title={Lösung},
  #1
}

% Glossar-Eintrag
\newcommand{\glossareintrag}[2]{%
  \par\noindent\textbf{\color{bfwPrimaryDark}#1} \enspace #2\par\smallskip
}

% Section-Stil — etwas farbiger als Standard
\usepackage{titlesec}
\titleformat{\section}
  {\Large\bfseries\color{bfwPrimaryDark}}
  {\thesection}{0.6em}{}
\titleformat{\subsection}
  {\large\bfseries\color{bfwPrimary}}
  {\thesubsection}{0.5em}{}
\titleformat{\subsubsection}
  {\normalsize\bfseries\color{bfwInk}}
  {\thesubsubsection}{0.4em}{}

% TikZ-Stil "skizzenhaft" — etwas weicher als Rough.js, aber stabil
\tikzset{
  roughbox/.style = {
    draw=bfwPrimaryDark, thick, fill=bfwBgTeal,
    rounded corners=4pt, inner sep=6pt
  },
  roughaccent/.style = {
    draw=bfwAccent, thick, fill=bfwBgAmber,
    rounded corners=4pt, inner sep=6pt
  },
  roughline/.style = {
    draw=bfwInkSoft, thick, -{Stealth[length=2.5mm]}
  }
}

% Bullet-Symbole (bewusst kein FontAwesome — vermeidet Font-Installations-Probleme)
\newcommand{\faLightbulb}{$\bullet$}
\newcommand{\faBookmark}{$\bullet$}

% Header/Footer
\usepackage{fancyhdr}
\pagestyle{fancy}
\fancyhf{}
\renewcommand{\headrulewidth}{0pt}
\fancyfoot[L]{\small\color{bfwInkSoft}\thepage}
\fancyfoot[R]{\small\color{bfwInkSoft} BFW M-064 Beschaffung im Büromanagement}


\title{\vspace*{-1.5cm}\Huge\bfseries\color{bfwPrimaryDark}Aufgabenheft Tag~11\\[0.4em]
  \large\color{bfwInkSoft}\textnormal{Postbearbeitung: Postausgang --- Übungen im IHK-Klausurformat}}
\author{\textsf{Büroprozesse gestalten und Arbeitsvorgänge organisieren \textperiodcentered{} M-067}\\
\textsf{\small\copyright\ Can Siebert\,\textperiodcentered\,2026}}
\date{Selbstlernmaterial \textperiodcentered{} 13.07.2026}

\begin{document}
\maketitle
\thispagestyle{empty}

\begin{merksatz}[title={\bfseries So nutzen Sie dieses Aufgabenheft}]
Bearbeiten Sie alle Aufgaben zunächst \emph{ohne} in die Lösungen zu schauen. Schreiben
Sie Ihre Antworten in vollständigen Sätzen --- die IHK-Prüfung verlangt formulierte
Begründungen, nicht bloße Stichworte. Beachten Sie die \emph{Operatoren} (nennen,
erläutern, beurteilen, begründen) --- sie geben den geforderten Antwort-Tiefegrad vor.
Erst danach öffnen Sie den Lösungsteil ab Seite~\pageref{loesungen} und vergleichen.
\end{merksatz}

\tableofcontents
\vspace{1em}
\hrule

\section{Aufgaben}

\subsection*{Aufgabe~1 --- Ablauf des Postausgangs (12 Punkte)}

\noindent\textbf{\color{bfwAccent}YouTube-Vorbereitung:}\enspace \href{https://www.youtube.com/watch?v=D6n_qJ4OOrs}{Postausgang --- Ablauf in der Poststelle} (\url{https://www.youtube.com/watch?v=D6n_qJ4OOrs})

\medskip
Sie arbeiten in der Poststelle der \emph{Lindner \& Partner GmbH} und bereiten die
Ausgangspost für den Versand vor.

\begin{enumerate}[label=(\alph*)]
  \item \textbf{Erläutern} Sie die Aufgabe der Poststelle beim Postausgang. \hfill (3~P.)
  \item \textbf{Bringen} Sie die sieben Arbeitsschritte in die richtige Reihenfolge. \hfill (4~P.)
  \item \textbf{Nennen} Sie drei Punkte, auf die beim Schritt „Kontrolle” geprüft wird. \hfill (3~P.)
  \item \textbf{Beurteilen} Sie, warum die feste Reihenfolge der Schritte sinnvoll ist. \hfill (2~P.)
\end{enumerate}

\subsection*{Aufgabe~2 --- Versand- und Sendungsarten (14 Punkte)}

\noindent\textbf{\color{bfwAccent}YouTube-Vorbereitung:}\enspace \href{https://www.youtube.com/watch?v=eEuG1GNZGOw}{Einschreiben-Arten der Deutschen Post} (\url{https://www.youtube.com/watch?v=eEuG1GNZGOw})

\medskip
Die Ausgangspost wird u.\,a.\ nach Versandart sortiert.

\begin{enumerate}[label=(\alph*)]
  \item \textbf{Erläutern} Sie den Unterschied zwischen \emph{Einschreiben Einwurf} und \emph{Einschreiben Übergabe}. \hfill (4~P.)
  \item \textbf{Erklären} Sie, wann sich \emph{Einschreiben Wert} anbietet. \hfill (3~P.)
  \item \textbf{Nennen} Sie Eignung und einen Nachteil der \emph{Warensendung}. \hfill (3~P.)
  \item \textbf{Beurteilen} Sie, welche Versandart sich für eine fristgebundene Kündigung eignet, und \textbf{begründen} Sie. \hfill (4~P.)
\end{enumerate}

\subsection*{Aufgabe~3 --- Falzen und Kuvertieren (12 Punkte)}

\noindent\textbf{\color{bfwAccent}YouTube-Vorbereitung:}\enspace \href{https://www.youtube.com/watch?v=gBOF5SJc3vg}{Falzarten --- Brief richtig falten} (\url{https://www.youtube.com/watch?v=gBOF5SJc3vg})

\medskip
\begin{enumerate}[label=(\alph*)]
  \item \textbf{Nennen} Sie vier übliche Falzarten. \hfill (4~P.)
  \item \textbf{Ordnen} Sie zu, welche Falzart einen A4-Briefbogen in eine Briefhülle DIN lang mit Sichtfenster bringt. \hfill (2~P.)
  \item \textbf{Erläutern} Sie, warum die Wahl der Falzart wirtschaftlich wichtig ist. \hfill (3~P.)
  \item \textbf{Beschreiben} Sie den Schritt „Kuvertieren und Verschließen”. \hfill (3~P.)
\end{enumerate}

\subsection*{Aufgabe~4 --- Frankieren und Aufgeben (15 Punkte)}

\noindent\textbf{\color{bfwAccent}YouTube-Vorbereitung:}\enspace \href{https://www.youtube.com/watch?v=eZGGxDmgxVg}{Frankiermaschine \& Frankierarten} (\url{https://www.youtube.com/watch?v=eZGGxDmgxVg})

\medskip
Die \emph{Lindner \& Partner GmbH} verschickt täglich rund 250 Briefe sowie monatlich
mehrere tausend Werbesendungen.

\begin{enumerate}[label=(\alph*)]
  \item \textbf{Nennen} Sie vier Frankierarten und \textbf{erläutern} Sie, welche für hohes Postaufkommen geeignet ist. \hfill (5~P.)
  \item \textbf{Erläutern} Sie den Unterschied zwischen Wertvorgabe- und Fernwertvorgabesystem. \hfill (4~P.)
  \item \textbf{Nennen} Sie drei Möglichkeiten, die Post aufzugeben. \hfill (3~P.)
  \item \textbf{Begründen} Sie, welchen Dienstleister Sie für ein eiliges, vertrauliches Vertragsoriginal empfehlen. \hfill (3~P.)
\end{enumerate}

\subsection*{Aufgabe~5 --- Poststraße und moderne Praxis (10 Punkte)}

\noindent\textbf{\color{bfwAccent}YouTube-Vorbereitung:}\enspace \href{https://www.youtube.com/watch?v=J6Im1kKtur4}{Poststraße \& Kuvertierstraße} (\url{https://www.youtube.com/watch?v=J6Im1kKtur4})

\medskip
\begin{enumerate}[label=(\alph*)]
  \item \textbf{Erklären} Sie, was eine Poststraße ist und welche Arbeitsschritte sie automatisch ausführt. \hfill (4~P.)
  \item \textbf{Erläutern} Sie, was moderne Versandsoftware leistet. \hfill (3~P.)
  \item \textbf{Erklären} Sie das Prinzip der Hybridpost. \hfill (3~P.)
\end{enumerate}

\subsection*{Aufgabe~6 --- Elektronische Post (12 Punkte)}

\noindent\textbf{\color{bfwAccent}YouTube-Vorbereitung:}\enspace \href{https://www.youtube.com/watch?v=ezaIWyAqonQ}{Pflichtangaben im Geschäftsbrief} (\url{https://www.youtube.com/watch?v=ezaIWyAqonQ})

\medskip
\begin{enumerate}[label=(\alph*)]
  \item \textbf{Erläutern} Sie, warum geschäftliche E-Mails Pflichtangaben enthalten müssen, und \textbf{nennen} Sie die gesetzliche Grundlage. \hfill (3~P.)
  \item \textbf{Nennen} Sie vier Pflichtangaben für den Geschäftsbrief einer GmbH. \hfill (4~P.)
  \item \textbf{Erläutern} Sie kurz De-Mail und EGVP. \hfill (3~P.)
  \item \textbf{Erklären} Sie die Funktion der elektronischen (qualifizierten) Signatur. \hfill (2~P.)
\end{enumerate}

\vspace{1em}
\noindent\textbf{Punkte gesamt: 75}

\newpage

\section{Lösungen}\label{loesungen}

\subsection*{Lösung Aufgabe~1}
\begin{loesung}
\textbf{(a)} In der Poststelle wird die in den einzelnen Abteilungen gesammelte
Ausgangspost zentral zusammengeführt und versandfertig gemacht. Dort fallen die einzelnen
Arbeitsschritte gebündelt an, was Zeit spart und einheitliche Qualität sichert.

\textbf{(b)} Reihenfolge: 1.~Kontrollieren, 2.~Sortieren, 3.~Falzen,
4.~Kuvertieren/Verschließen, 5.~Wiegen, 6.~Frankieren, 7.~Aufgeben.

\textbf{(c)} Geprüft wird auf \emph{Unterschrift}, \emph{Adresse} und \emph{Vollständigkeit
der Anlagen}. Fehlt bei einem Nicht-Fensterumschlag die Adresse, wird sie per
Adressiermaschine oder Adressaufkleber ergänzt.

\textbf{(d)} Die Schritte bauen aufeinander auf: Erst nach Kontrolle und Sortierung steht
fest, welches Format und welche Versandart nötig sind; gewogen und frankiert werden kann
erst, wenn der Brief gefalzt und kuvertiert ist. Eine falsche Reihenfolge führt zu
Fehlern und unnötigen Portokosten.
\end{loesung}

\subsection*{Lösung Aufgabe~2}
\begin{loesung}
\textbf{(a)} Beim \emph{Einschreiben Einwurf} dokumentiert der Postzusteller lediglich,
dass der Brief in den Briefkasten oder das Postfach des Empfängers eingeworfen wurde.
Beim \emph{Einschreiben (Übergabe)} wird der Brief gegen Unterschrift an den Empfänger
oder einen Empfangsberechtigten übergeben; eine Sendungsverfolgung im Internet ist
möglich.

\textbf{(b)} \emph{Einschreiben Wert} bietet sich an, wenn Bargeld (bis 100,00~€) oder
kleinere Sachwerte (bis 500,00~€) per Brief verschickt werden. Der Versender erhält einen
Einlieferungsnachweis, kann die Sendung verfolgen, und die Zustellung wird dokumentiert.

\textbf{(c)} Die \emph{Warensendung} eignet sich für kleinere Gegenstände und
Werbesendungen (z.\,B.\ Kataloge, Ersatzteile) und ist kostengünstig. Nachteil: Sie ist
nicht versichert, und eine Sendungsverfolgung ist nicht möglich.

\textbf{(d)} Für eine fristgebundene Kündigung eignet sich das \emph{Einschreiben
Rückschein} (oder Einschreiben Einwurf). Begründung: Die dokumentierte Zustellung
beweist, dass und wann der Empfänger den Brief erhalten hat --- wichtig bei
fristgebundenen, rechtserheblichen Schriftstücken.
\end{loesung}

\subsection*{Lösung Aufgabe~3}
\begin{loesung}
\textbf{(a)} Vier Falzarten: Bruch-/Einfachfalz, Kreuzfalz, Zickzack-/Leporellofalz,
Wickelfalz.

\textbf{(b)} Der \emph{Zickzack- bzw.\ Leporellofalz} bringt einen A4-Briefbogen (auf
1/3~A4 gefaltet) in eine Briefhülle DIN lang mit Sichtfenster; auch der Wickelfalz führt
zu DIN lang.

\textbf{(c)} Die richtige Falzart bestimmt das Format und damit das Porto. Ein möglichst
kleines, zulässiges Format senkt die Versandkosten; ein zu großes Format verteuert den
Versand unnötig --- bei hohem Postaufkommen ein erheblicher Kostenfaktor.

\textbf{(d)} Beim \emph{Kuvertieren} werden die Schriftstücke in die für den Versand
geeigneten Umschläge gesteckt --- manuell oder maschinell mit Kuvertiermaschinen.
Anschließend werden die Briefe \emph{verschlossen}.
\end{loesung}

\subsection*{Lösung Aufgabe~4}
\begin{loesung}
\textbf{(a)} Vier Frankierarten: Briefmarken (Postwertzeichen), Plusbrief (Umschlag mit
aufgedruckter Briefmarke), DV-Verfahren (Frankieren mittels EDV, genehmigungspflichtig),
Frankiermaschine. Für hohes Postaufkommen eignen sich \emph{DV-Verfahren} und
\emph{Frankiermaschine} (auch Online-Frankierung), da sie Massenpost schnell und ohne
einzelnes Aufkleben von Marken freimachen.

\textbf{(b)} Beim \emph{Wertvorgabesystem} wird ein Betrag im Voraus bezahlt, in der
Frankiermaschine eingestellt und durch das Stempeln verbraucht. Beim
\emph{Fernwertvorgabesystem} wird der Entgelt-Betrag per Telefon oder PC beantragt und
übermittelt und dann mithilfe eines Codes in die Frankiermaschine eingegeben.

\textbf{(c)} Aufgeben in einer \emph{Postfiliale}, über den \emph{Bring- und
Abholservice} (z.\,B.\ HIN + WEG) oder über \emph{KEP-Dienste} (Kurier-, Express-,
Paketdienste).

\textbf{(d)} Empfehlung: \emph{Kurierdienst}. Er befördert die Sendung persönlich und
direkt zum Empfänger, ist schnell (z.\,B.\ Fahrradkuriere in Großstädten) und eignet sich
für vertrauliche, besonders schutzbedürftige Dokumente. Soll zusätzlich eine feste Frist
garantiert sein, kommt ein \emph{Expressdienst} in Frage.
\end{loesung}

\subsection*{Lösung Aufgabe~5}
\begin{loesung}
\textbf{(a)} Eine \emph{Poststraße} entsteht, wenn die für die Arbeitsschritte benötigten
Maschinen hintereinandergeschaltet und in einem Raum zentral zusammengefasst werden. Sie
ermöglicht \emph{Falzen, Kuvertieren, Verschließen, Trennen nach Portoklassen und
Frankieren} in einem automatischen Maschinengang.

\textbf{(b)} \emph{Versandsoftware} ermittelt automatisch die günstigste Versandart und
das Porto, erstellt Labels und Frankierungen, verwaltet Adressen und übergibt die Daten
elektronisch an den Dienstleister.

\textbf{(c)} Bei der \emph{Hybridpost} übergibt das Unternehmen die Briefe nur noch
digital; ein Dienstleister druckt, kuvertiert, frankiert und stellt sie physisch zu. Der
Medienbruch von digital zu Papier erfolgt erst beim Dienstleister.
\end{loesung}

\subsection*{Lösung Aufgabe~6}
\begin{loesung}
\textbf{(a)} Geschäftliche E-Mails ersetzen den postalischen Geschäftsbrief; daher gelten
für sie dieselben Pflichtangaben. Gesetzliche Grundlage ist das \emph{EHUG} (Gesetz über
das elektronische Handels-, Genossenschafts- und Unternehmensregister) bzw.\ die
EU-Publizitätsrichtlinie: Die Pflichtangaben gelten unabhängig von der Form (Papier oder
elektronisch).

\textbf{(b)} Vier Pflichtangaben einer GmbH: Firma, Rechtsform (Abkürzung GmbH genügt),
Sitz der Gesellschaft, Registergericht und Handelsregisternummer (sowie alle
Geschäftsführer).

\textbf{(c)} \emph{De-Mail} ist ein verschlüsselter, nachweisbarer E-Mail-Dienst für
verbindliche elektronische Kommunikation. Das \emph{EGVP} (Elektronisches Gerichts- und
Verwaltungspostfach) dient dem rechtsverbindlichen elektronischen Schriftverkehr mit
Gerichten und Behörden.

\textbf{(d)} Die qualifizierte elektronische Signatur ersetzt die eigenhändige
Unterschrift. Sie weist die Identität des Absenders nach und stellt sicher, dass das
Dokument nach dem Signieren nicht verändert wurde (Integrität) --- so entsteht
Rechtssicherheit im digitalen Versand.
\end{loesung}

\end{document}
